% =========================================================
% LMTH 2040 — Fall 2026
% Core Reader Plan
%
% This file is designed to be included from syllabus.tex with:
%     % =========================================================
% LMTH 2040 — Fall 2026
% Core Reader Plan
%
% This file is designed to be included from syllabus.tex with:
%     % =========================================================
% LMTH 2040 — Fall 2026
% Core Reader Plan
%
% This file is designed to be included from syllabus.tex with:
%     % =========================================================
% LMTH 2040 — Fall 2026
% Core Reader Plan
%
% This file is designed to be included from syllabus.tex with:
%     \input{readings}
%
% Citations assume references.bib is loaded in syllabus.tex via:
%     \addbibresource{references.bib}
%
% Page ranges below are working selections. Verify against the
% exact editions placed on Canvas / Course Reserves before release.
% =========================================================

\section{Reader}

The readings in this course are not supplementary background to the
mathematics. They are part of our investigation of what calculus is, how
mathematical knowledge develops, how models and measurements acquire authority,
and how ideas of change, optimization, representation, and intelligence have
been transformed across different historical settings.

The reader moves among three kinds of texts:

\begin{itemize}[leftmargin=*]
    \item \textbf{Primary sources}: mathematical, scientific, and technical
    texts written by historical actors as ideas were being developed;
    \item \textbf{History and sociology of science}: scholarship examining
    mathematics, models, objectivity, quantification, and scientific practice
    historically and socially;
    \item \textbf{Theory and cultural studies}: texts that help us question
    space, classification, representation, images, technology, and artificial
    intelligence.
\end{itemize}

Most reading assignments are deliberately short. Read slowly and bring the
text into contact with the mathematics. For each encounter, use the three
questions below:

\begin{description}[leftmargin=1.15in,style=nextline]
    \item[\textcolor{LangRed}{READ}] What is the author actually claiming,
    constructing, or trying to accomplish?
    \item[\textcolor{LangRed}{CONNECT}] What happens when we place the text
    beside the mathematics, computation, model, or image we are studying?
    \item[\textcolor{LangRed}{QUESTION}] Where should we challenge, test,
    complicate, or refuse the author's account?
\end{description}

% ---------------------------------------------------------
\subsection*{Encounter 1 \hfill August 26}
\courselabel{Genesis: What is calculus?}

\textbf{Otto Toeplitz, \emph{The Calculus: A Genetic Approach}.}
Introductory selection. \parencite{toeplitz1963}

\textit{Type: History of mathematics / pedagogy}

\begin{description}[leftmargin=1.15in,style=nextline]
\item[\textcolor{LangRed}{READ}]
What does Toeplitz mean by a ``genetic'' approach to mathematics?
\item[\textcolor{LangRed}{CONNECT}]
Compare the order in which mathematical ideas historically arise with the
order in which a modern calculus textbook usually presents them.
\item[\textcolor{LangRed}{QUESTION}]
Should the historical genesis of a mathematical idea determine how we learn
it? Where might historical order help, and where might it mislead?
\end{description}

% ---------------------------------------------------------
\subsection*{Encounter 2 \hfill September 2}
\courselabel{Archimedes: Measurement, exhaustion, and matter}

\textbf{Archimedes, \emph{Quadrature of the Parabola}.}
Short primary-source selection from the geometric argument.

\textit{Type: Primary mathematical source}

\medskip

\textbf{Michel Serres, \emph{The Birth of Physics}.}
``The Work of Archimedes,'' working selection pp.~32--36, with a short
selection from ``Archimedes on Thinking Deviation.''

\textit{Type: Philosophy / history of science}

\begin{description}[leftmargin=1.15in,style=nextline]
\item[\textcolor{LangRed}{READ}]
How does Archimedes reason about an area without possessing our modern
definition of the definite integral? What does Serres think is at stake in
Archimedean mathematics?
\item[\textcolor{LangRed}{CONNECT}]
Identify the finite approximations in Archimedes' argument. How are they like,
and unlike, the computational approximations we construct in Jupyter?
\item[\textcolor{LangRed}{QUESTION}]
Does describing Archimedes as someone who ``almost discovered calculus'' help
us understand his work, or erase the mathematical world in which he worked?
\end{description}

% ---------------------------------------------------------
% ---------------------------------------------------------
\subsection*{Encounter 3 \hfill September 14}
\courselabel{Indivisibles, the continuum, and nomad science}

\textbf{Bonaventura Cavalieri, ``Principle of Cavalieri'' and
``Integration,'' in \textcite{struik1986}, selections beginning
pp.~209 and 214.}

\textit{Type: Primary mathematical source}

\medskip

\textbf{Gilles Deleuze and Félix Guattari, ``1227: Treatise on
Nomadology---The War Machine,'' in \emph{A Thousand Plateaus}.}

Selected passages on nomad science and royal or State science.
\parencite{deleuzeguattari1987}

\textit{Type: Philosophy / cultural theory}

\begin{description}[leftmargin=1.15in,style=nextline]

\item[\textcolor{LangRed}{READ}]
How does Cavalieri reason with indivisibles? What distinction do
Deleuze and Guattari draw between nomad and royal forms of knowledge?

\item[\textcolor{LangRed}{CONNECT}]
Does Cavalieri's treatment of continuous magnitudes fit Deleuze and
Guattari's description of a mathematical practice concerned with
becoming, variation, and local operations rather than fixed forms?

\item[\textcolor{LangRed}{QUESTION}]
Does ``nomad science'' actually illuminate Cavalieri's mathematics,
or are we imposing a twentieth-century philosophical category on a
seventeenth-century practice?

\end{description}

% ---------------------------------------------------------
\subsection*{Encounter 4 \hfill September 16}
\courselabel{Smooth and striated space}

\textbf{Deleuze and Guattari, \emph{A Thousand Plateaus}.}
Short selected passage on smooth and striated space.

\textit{Type: Philosophy / cultural theory}

\begin{description}[leftmargin=1.15in,style=nextline]
\item[\textcolor{LangRed}{READ}]
How do smooth and striated spaces differ in the authors' account?
\item[\textcolor{LangRed}{CONNECT}]
A Riemann sum imposes a coordinate system, partitions an interval, samples a
function, assigns rectangles, and aggregates measurements. In what sense does
this ``striate'' a continuous object?
\item[\textcolor{LangRed}{QUESTION}]
What becomes mathematically possible because of the partition? What information
or experience of the continuous object might the partition suppress?
\end{description}

% ---------------------------------------------------------
\subsection*{Encounter 5 \hfill September 23}
\courselabel{Before the Fundamental Theorem was fundamental}

\textbf{Isaac Barrow, ``The Fundamental Theorem of the Calculus,''
in \textcite{struik1986}, selection beginning p.~253.}

\textit{Type: Primary mathematical source}

\begin{description}[leftmargin=1.15in,style=nextline]

\item[\textcolor{LangRed}{READ}]
What relationship does Barrow establish between problems involving
tangents and problems involving areas?

\item[\textcolor{LangRed}{CONNECT}]
Translate Barrow's geometric reasoning into our modern relationship
between accumulation and differentiation:

\[
F(x)=\int_a^x f(t)\,dt
\qquad\Longrightarrow\qquad
F'(x)=f(x).
\]

What does our notation make easier to see?

\item[\textcolor{LangRed}{QUESTION}]
What gets lost when we call Barrow's argument ``the Fundamental
Theorem of Calculus'' before the modern objects called derivatives
and integrals have been stabilized?

\end{description}

% ---------------------------------------------------------
\subsection*{Encounter 6 \hfill September 30}
\courselabel{Quantification: When inequality becomes a number}

\textbf{Theodore M. Porter, \emph{Trust in Numbers}.}
Working selection from ``Cultures of Objectivity'' and ``How Social Numbers
Are Made Valid.'' \parencite{porter1995}

\textit{Type: History / sociology of quantification}

\begin{description}[leftmargin=1.15in,style=nextline]
\item[\textcolor{LangRed}{READ}]
Why, for Porter, can quantitative procedures acquire authority precisely by
appearing to reduce personal judgment?
\item[\textcolor{LangRed}{CONNECT}]
Construct a Lorenz curve and calculate a Gini coefficient. Identify at least
three decisions required before ``inequality'' can be represented as an area
and then as a single number.
\item[\textcolor{LangRed}{QUESTION}]
Which decisions in the Gini calculation are mathematical and which are social?
Is that distinction as clear as it first appears?
\end{description}

% ---------------------------------------------------------


% ---------------------------------------------------------
% ---------------------------------------------------------
\subsection*{Encounter 7 \hfill October 14}
\courselabel{The differential calculus appears}

\textbf{Gottfried Wilhelm Leibniz, ``The First Publication of His
Differential Calculus,'' in \textcite{struik1986},
selection beginning p.~271.}

\textit{Type: Primary mathematical source}

\begin{description}[leftmargin=1.15in,style=nextline]

\item[\textcolor{LangRed}{READ}]
What mathematical objects and operations does Leibniz introduce?
How does his notation participate in the argument?

\item[\textcolor{LangRed}{CONNECT}]
Translate a portion of Leibniz's argument into contemporary notation.
Identify where we would now speak about a derivative, differential,
product rule, quotient rule, or tangent.

\item[\textcolor{LangRed}{QUESTION}]
Does translating Leibniz into modern notation help us understand him,
or does it make his mathematics look more like ours than it really was?

\end{description}

% ---------------------------------------------------------
\subsection*{Encounter 8 \hfill October 26}
\courselabel{Seeing change: Images and objectivity}

\textbf{Lorraine Daston and Peter Galison,
\emph{Objectivity}.}

Selected passages on mechanical objectivity and scientific images.
\parencite{dastongalison2007}

\textit{Type: History of science / visual culture}

\medskip

\textbf{Visual archive: selected motion studies by
Eadweard Muybridge and Étienne-Jules Marey.}

\textit{Type: Primary visual sources}

\begin{description}[leftmargin=1.15in,style=nextline]

\item[\textcolor{LangRed}{READ}]
What does ``mechanical objectivity'' promise? What practices of
seeing does it replace or reorganize?

\item[\textcolor{LangRed}{CONNECT}]
Compare a mathematical graph of motion, a sequence of photographs,
a table of measurements, and a numerical derivative. What form of
change does each make visible?

\item[\textcolor{LangRed}{QUESTION}]
Is an image evidence, measurement, interpretation, or model?
Can it be all four?

\end{description}

% ---------------------------------------------------------
\subsection*{Encounter 9 \hfill November 4}
\courselabel{Slope, extrema, and optimization}

\textbf{Michel Serres, \emph{The Birth of Physics}.}

Selected passage from ``Slope and Extrema,'' beginning around p.~52.

\textit{Type: Philosophy / history of science}

\begin{description}[leftmargin=1.15in,style=nextline]

\item[\textcolor{LangRed}{READ}]
How does Serres connect slope, deviation, equilibrium, and extrema
to a physical picture of the world?

\item[\textcolor{LangRed}{CONNECT}]
If a function is a landscape, its derivative tells us which way
the landscape slopes. An extremum is a place where that local
direction vanishes.

What changes when we turn this geometric fact into an algorithm:
repeatedly moving in the direction indicated by the derivative?

\item[\textcolor{LangRed}{QUESTION}]
Is optimization simply a mathematical technique, or does it encourage
a particular way of imagining a problem---as a landscape with states,
directions, extrema, and preferred paths?

\end{description}


% ---------------------------------------------------------
\subsection*{Encounter 8 \hfill November 9}
\courselabel{What does it mean for a machine to learn?}

\textbf{Frank Rosenblatt, ``The Perceptron: A Probabilistic Model for
Information Storage and Organization in the Brain.''}
Selected opening and model passages from the 1958 paper.

\textit{Type: Primary source / history of artificial intelligence}

\begin{description}[leftmargin=1.15in,style=nextline]
\item[\textcolor{LangRed}{READ}]
What does Rosenblatt claim the perceptron can model, and what does he mean by
learning or self-organization?
\item[\textcolor{LangRed}{CONNECT}]
In our perceptron model, identify the inputs, weights, weighted sum, decision
rule, error, and update. Exactly what changes when the machine ``learns''?
\item[\textcolor{LangRed}{QUESTION}]
Formulate the strongest mathematical argument you can for calling parameter
adjustment ``learning,'' and the strongest conceptual argument against doing so.
\end{description}

% ---------------------------------------------------------
\subsection*{Encounter 9 \hfill November 11}
\courselabel{The chain rule becomes a learning procedure}

\textbf{David E. Rumelhart, Geoffrey E. Hinton, and Ronald J. Williams,
``Learning Representations by Back-Propagating Errors.''}
\emph{Nature} 323 (1986), pp.~533--536.

\textit{Type: Primary technical source / artificial intelligence}

\begin{description}[leftmargin=1.15in,style=nextline]
\item[\textcolor{LangRed}{READ}]
What problem does backpropagation solve, and what do the authors claim hidden
units learn to represent?
\item[\textcolor{LangRed}{CONNECT}]
Annotate the paper using the mathematics of the course: Where is the function?
Where is composition? Where is the loss? Where is the derivative? Where is the
chain rule? Where is optimization?
\item[\textcolor{LangRed}{QUESTION}]
What has to happen conceptually before a sequence of derivative calculations
can be described as ``learning representations''?
\end{description}

% ---------------------------------------------------------
\subsection*{Encounter 10 \hfill November 16}
\courselabel{Flows and paths}

\textbf{Michel Serres, \emph{The Birth of Physics}.}
Working selection from ``Flows and Paths,'' beginning around p.~70.

\textit{Type: Philosophy / history of science}

\begin{description}[leftmargin=1.15in,style=nextline]
\item[\textcolor{LangRed}{READ}]
What distinction does Serres make between describing an object and describing
a flow, path, or process?
\item[\textcolor{LangRed}{CONNECT}]
A differential equation gives a rule for change rather than directly giving a
trajectory. How does this alter what it means to know a mathematical object?
\item[\textcolor{LangRed}{QUESTION}]
Does the language of flow illuminate differential equations, or become merely
metaphorical once we formalize the mathematics?
\end{description}

% ---------------------------------------------------------
\subsection*{Encounter 11 \hfill November 18}
\courselabel{Can a model be verified?}

\textbf{Naomi Oreskes, Kristin Shrader-Frechette, and Kenneth Belitz,
``Verification, Validation, and Confirmation of Numerical Models in the Earth
Sciences.''}
\emph{Science} 263 (1994), pp.~641--646.

\textit{Type: Philosophy / history of modeling}

\begin{description}[leftmargin=1.15in,style=nextline]
\item[\textcolor{LangRed}{READ}]
Why do the authors distinguish confirmation from verification or validation for
models of open natural systems?
\item[\textcolor{LangRed}{CONNECT}]
Return to an Euler-method simulation. Which parts come from mathematics, which
from empirical assumptions, which from numerical choices, and which from
initial conditions?
\item[\textcolor{LangRed}{QUESTION}]
If a model produces accurate predictions, what has actually been established?
What has not?
\end{description}

% ---------------------------------------------------------
\subsection*{Encounter 12 \hfill November 30}
\courselabel{Feedback, control, and intelligence}

\textbf{Norbert Wiener, \emph{The Human Use of Human Beings}.}
Short selection on cybernetics, feedback, communication, and purposive behavior.
\parencite{wiener1950}

\textit{Type: Primary source / cybernetics}

\medskip

\textbf{Andrew Pickering, \emph{The Cybernetic Brain}.}
Short historical selection on cybernetic practice.

\textit{Type: History / sociology of science}

\begin{description}[leftmargin=1.15in,style=nextline]
\item[\textcolor{LangRed}{READ}]
How does feedback permit Wiener and other cyberneticians to describe purposive
behavior without beginning from an internal mind or intention?
\item[\textcolor{LangRed}{CONNECT}]
Identify feedback mathematically in one of our dynamical systems. How do
equilibrium and stability change when the output of a system becomes part of
its future input?
\item[\textcolor{LangRed}{QUESTION}]
What is gained---and what may be lost---when intelligence, behavior, or purpose
is redescribed in terms of information and feedback?
\end{description}

% ---------------------------------------------------------
% ---------------------------------------------------------
\subsection*{Encounter 13 \hfill December 2}
\courselabel{When an equation becomes a field}

\textbf{Jean le Rond d'Alembert, Leonhard Euler, and Daniel Bernoulli,
``The Vibrating String and Its Partial Differential Equation,''
in \textcite{struik1986}, selection beginning p.~351.}

\textit{Type: Primary mathematical / physical source}

\begin{description}[leftmargin=1.15in,style=nextline]

\item[\textcolor{LangRed}{READ}]
What is being represented mathematically when the vibrating string is
described not by a single changing quantity but by a function of both
position and time?

\item[\textcolor{LangRed}{CONNECT}]
Compare

\[
y=y(t)
\]

with

\[
u=u(x,t).
\]

Why does the second description require partial derivatives? What
does each independent variable represent?

\item[\textcolor{LangRed}{QUESTION}]
Is the partial differential equation a description of the string,
a model of the string, or a new mathematical object with a life
independent of the physical system that produced it?

\end{description}



% ---------------------------------------------------------
\subsection*{Encounter 14 \hfill December 7}
\courselabel{Surfaces, dimensions, and classification}

\textbf{Gaspard Monge, ``The Two Curvatures of a Curved Surface,''
in \textcite{struik1986}, short selection beginning p.~413.}

\textit{Type: Primary mathematical source}

\medskip

\textbf{Geoffrey C. Bowker and Susan Leigh Star,
\emph{Sorting Things Out: Classification and Its Consequences}.}

Selected introduction. \parencite{bowkerstar1999}

\textit{Type: Science and technology studies}

\begin{description}[leftmargin=1.15in,style=nextline]

\item[\textcolor{LangRed}{READ}]
How does moving from a curve to a surface change the mathematical
problem? How do Bowker and Star describe the work performed by
systems of classification?

\item[\textcolor{LangRed}{CONNECT}]
For a curve

\[
y=f(x),
\]

local change can be described by a single derivative. For a surface

\[
z=f(x,y),
\]

change depends upon direction.

How does moving from curves to surfaces change the meaning of slope,
tangent, and curvature? How does this prepare us to think about
high-dimensional spaces in which objects are represented by many
coordinates?

\item[\textcolor{LangRed}{QUESTION}]
An embedding represents an object as a point in a high-dimensional
space. What decisions must occur before similarity, difference, and
classification can appear as geometric relationships?

\end{description}

% ---------------------------------------------------------
\subsection*{Encounter 17 \hfill December 9}
\courselabel{Images for machines}

\textbf{Trevor Paglen, ``Invisible Images (Your Pictures Are Looking at
You).''}

\textit{Type: Art / visual culture / critical technology studies}

\begin{description}[leftmargin=1.15in,style=nextline]
\item[\textcolor{LangRed}{READ}]
What changes, for Paglen, when images circulate primarily between machines
rather than between images and human viewers?
\item[\textcolor{LangRed}{CONNECT}]
Trace the transformations
\[
\text{image} \longrightarrow \text{pixels} \longrightarrow
\text{numbers} \longrightarrow \text{features} \longrightarrow
\text{vector / embedding} \longrightarrow \text{classification}.
\]
At which stages is calculus involved?
\item[\textcolor{LangRed}{QUESTION}]
If no human needs to see an image for the image to act in the world, should we
still understand it primarily as a visual object?
\end{description}

% ---------------------------------------------------------
\subsection*{Encounter 18 \hfill December 14}
\courselabel{Generative AI: Mathematics, infrastructure, and culture}

\textbf{Kate Crawford, \emph{Atlas of AI}.}
Selected concluding or synthetic passage. \parencite{crawford2021}

\textit{Type: Critical AI / sociology of technology}

\medskip

\textbf{Contemporary primary technical artifact.}
A short excerpt, diagram, or model description from work on diffusion or
generative models, to be selected.

\textit{Type: Primary technical source}

\begin{description}[leftmargin=1.15in,style=nextline]
\item[\textcolor{LangRed}{READ}]
What material, institutional, political, and labor systems does Crawford ask us
to see when we use the term ``artificial intelligence''?
\item[\textcolor{LangRed}{CONNECT}]
Our mathematical story may describe a generative system through probability,
gradients, fields, differential equations, or optimization. What does that
description explain extremely well? What does it leave outside the frame?
\item[\textcolor{LangRed}{QUESTION}]
What would an adequate account of generative AI have to contain in addition to
its mathematics?
\end{description}

% =========================================================
\section{Calculus Colloquium: Extended Reader}

The shared reader is intentionally small. The Calculus Colloquium expands it.
Students will select additional primary sources, historical scholarship,
theoretical texts, mathematical problems, visual artifacts, datasets, or
technical papers related to the current module.

Possible directions include Archimedes, Lucretius, Cavalieri, Newton, Leibniz,
Euler, Fourier, Marey, Muybridge, Turing, Shannon, McCulloch and Pitts, Hebb,
Rosenblatt, Wiener, Minsky and Papert, Rumelhart--Hinton--Williams, Hopfield,
convolutional networks, word embeddings, attention, transformers, GANs,
diffusion models, cybernetic art, Project Cybersyn, ImageNet, algorithmic
classification, scientific visualization, and histories of computation.

A colloquium contribution should not simply summarize its source. It should
make an argument about the relationship between the source and the mathematics
of the course.

%
% Citations assume references.bib is loaded in syllabus.tex via:
%     \addbibresource{references.bib}
%
% Page ranges below are working selections. Verify against the
% exact editions placed on Canvas / Course Reserves before release.
% =========================================================

\section{Reader}

The readings in this course are not supplementary background to the
mathematics. They are part of our investigation of what calculus is, how
mathematical knowledge develops, how models and measurements acquire authority,
and how ideas of change, optimization, representation, and intelligence have
been transformed across different historical settings.

The reader moves among three kinds of texts:

\begin{itemize}[leftmargin=*]
    \item \textbf{Primary sources}: mathematical, scientific, and technical
    texts written by historical actors as ideas were being developed;
    \item \textbf{History and sociology of science}: scholarship examining
    mathematics, models, objectivity, quantification, and scientific practice
    historically and socially;
    \item \textbf{Theory and cultural studies}: texts that help us question
    space, classification, representation, images, technology, and artificial
    intelligence.
\end{itemize}

Most reading assignments are deliberately short. Read slowly and bring the
text into contact with the mathematics. For each encounter, use the three
questions below:

\begin{description}[leftmargin=1.15in,style=nextline]
    \item[\textcolor{LangRed}{READ}] What is the author actually claiming,
    constructing, or trying to accomplish?
    \item[\textcolor{LangRed}{CONNECT}] What happens when we place the text
    beside the mathematics, computation, model, or image we are studying?
    \item[\textcolor{LangRed}{QUESTION}] Where should we challenge, test,
    complicate, or refuse the author's account?
\end{description}

% ---------------------------------------------------------
\subsection*{Encounter 1 \hfill August 26}
\courselabel{Genesis: What is calculus?}

\textbf{Otto Toeplitz, \emph{The Calculus: A Genetic Approach}.}
Introductory selection. \parencite{toeplitz1963}

\textit{Type: History of mathematics / pedagogy}

\begin{description}[leftmargin=1.15in,style=nextline]
\item[\textcolor{LangRed}{READ}]
What does Toeplitz mean by a ``genetic'' approach to mathematics?
\item[\textcolor{LangRed}{CONNECT}]
Compare the order in which mathematical ideas historically arise with the
order in which a modern calculus textbook usually presents them.
\item[\textcolor{LangRed}{QUESTION}]
Should the historical genesis of a mathematical idea determine how we learn
it? Where might historical order help, and where might it mislead?
\end{description}

% ---------------------------------------------------------
\subsection*{Encounter 2 \hfill September 2}
\courselabel{Archimedes: Measurement, exhaustion, and matter}

\textbf{Archimedes, \emph{Quadrature of the Parabola}.}
Short primary-source selection from the geometric argument.

\textit{Type: Primary mathematical source}

\medskip

\textbf{Michel Serres, \emph{The Birth of Physics}.}
``The Work of Archimedes,'' working selection pp.~32--36, with a short
selection from ``Archimedes on Thinking Deviation.''

\textit{Type: Philosophy / history of science}

\begin{description}[leftmargin=1.15in,style=nextline]
\item[\textcolor{LangRed}{READ}]
How does Archimedes reason about an area without possessing our modern
definition of the definite integral? What does Serres think is at stake in
Archimedean mathematics?
\item[\textcolor{LangRed}{CONNECT}]
Identify the finite approximations in Archimedes' argument. How are they like,
and unlike, the computational approximations we construct in Jupyter?
\item[\textcolor{LangRed}{QUESTION}]
Does describing Archimedes as someone who ``almost discovered calculus'' help
us understand his work, or erase the mathematical world in which he worked?
\end{description}

% ---------------------------------------------------------
% ---------------------------------------------------------
\subsection*{Encounter 3 \hfill September 14}
\courselabel{Indivisibles, the continuum, and nomad science}

\textbf{Bonaventura Cavalieri, ``Principle of Cavalieri'' and
``Integration,'' in \textcite{struik1986}, selections beginning
pp.~209 and 214.}

\textit{Type: Primary mathematical source}

\medskip

\textbf{Gilles Deleuze and Félix Guattari, ``1227: Treatise on
Nomadology---The War Machine,'' in \emph{A Thousand Plateaus}.}

Selected passages on nomad science and royal or State science.
\parencite{deleuzeguattari1987}

\textit{Type: Philosophy / cultural theory}

\begin{description}[leftmargin=1.15in,style=nextline]

\item[\textcolor{LangRed}{READ}]
How does Cavalieri reason with indivisibles? What distinction do
Deleuze and Guattari draw between nomad and royal forms of knowledge?

\item[\textcolor{LangRed}{CONNECT}]
Does Cavalieri's treatment of continuous magnitudes fit Deleuze and
Guattari's description of a mathematical practice concerned with
becoming, variation, and local operations rather than fixed forms?

\item[\textcolor{LangRed}{QUESTION}]
Does ``nomad science'' actually illuminate Cavalieri's mathematics,
or are we imposing a twentieth-century philosophical category on a
seventeenth-century practice?

\end{description}

% ---------------------------------------------------------
\subsection*{Encounter 4 \hfill September 16}
\courselabel{Smooth and striated space}

\textbf{Deleuze and Guattari, \emph{A Thousand Plateaus}.}
Short selected passage on smooth and striated space.

\textit{Type: Philosophy / cultural theory}

\begin{description}[leftmargin=1.15in,style=nextline]
\item[\textcolor{LangRed}{READ}]
How do smooth and striated spaces differ in the authors' account?
\item[\textcolor{LangRed}{CONNECT}]
A Riemann sum imposes a coordinate system, partitions an interval, samples a
function, assigns rectangles, and aggregates measurements. In what sense does
this ``striate'' a continuous object?
\item[\textcolor{LangRed}{QUESTION}]
What becomes mathematically possible because of the partition? What information
or experience of the continuous object might the partition suppress?
\end{description}

% ---------------------------------------------------------
\subsection*{Encounter 5 \hfill September 23}
\courselabel{Before the Fundamental Theorem was fundamental}

\textbf{Isaac Barrow, ``The Fundamental Theorem of the Calculus,''
in \textcite{struik1986}, selection beginning p.~253.}

\textit{Type: Primary mathematical source}

\begin{description}[leftmargin=1.15in,style=nextline]

\item[\textcolor{LangRed}{READ}]
What relationship does Barrow establish between problems involving
tangents and problems involving areas?

\item[\textcolor{LangRed}{CONNECT}]
Translate Barrow's geometric reasoning into our modern relationship
between accumulation and differentiation:

\[
F(x)=\int_a^x f(t)\,dt
\qquad\Longrightarrow\qquad
F'(x)=f(x).
\]

What does our notation make easier to see?

\item[\textcolor{LangRed}{QUESTION}]
What gets lost when we call Barrow's argument ``the Fundamental
Theorem of Calculus'' before the modern objects called derivatives
and integrals have been stabilized?

\end{description}

% ---------------------------------------------------------
\subsection*{Encounter 6 \hfill September 30}
\courselabel{Quantification: When inequality becomes a number}

\textbf{Theodore M. Porter, \emph{Trust in Numbers}.}
Working selection from ``Cultures of Objectivity'' and ``How Social Numbers
Are Made Valid.'' \parencite{porter1995}

\textit{Type: History / sociology of quantification}

\begin{description}[leftmargin=1.15in,style=nextline]
\item[\textcolor{LangRed}{READ}]
Why, for Porter, can quantitative procedures acquire authority precisely by
appearing to reduce personal judgment?
\item[\textcolor{LangRed}{CONNECT}]
Construct a Lorenz curve and calculate a Gini coefficient. Identify at least
three decisions required before ``inequality'' can be represented as an area
and then as a single number.
\item[\textcolor{LangRed}{QUESTION}]
Which decisions in the Gini calculation are mathematical and which are social?
Is that distinction as clear as it first appears?
\end{description}

% ---------------------------------------------------------


% ---------------------------------------------------------
% ---------------------------------------------------------
\subsection*{Encounter 7 \hfill October 14}
\courselabel{The differential calculus appears}

\textbf{Gottfried Wilhelm Leibniz, ``The First Publication of His
Differential Calculus,'' in \textcite{struik1986},
selection beginning p.~271.}

\textit{Type: Primary mathematical source}

\begin{description}[leftmargin=1.15in,style=nextline]

\item[\textcolor{LangRed}{READ}]
What mathematical objects and operations does Leibniz introduce?
How does his notation participate in the argument?

\item[\textcolor{LangRed}{CONNECT}]
Translate a portion of Leibniz's argument into contemporary notation.
Identify where we would now speak about a derivative, differential,
product rule, quotient rule, or tangent.

\item[\textcolor{LangRed}{QUESTION}]
Does translating Leibniz into modern notation help us understand him,
or does it make his mathematics look more like ours than it really was?

\end{description}

% ---------------------------------------------------------
\subsection*{Encounter 8 \hfill October 26}
\courselabel{Seeing change: Images and objectivity}

\textbf{Lorraine Daston and Peter Galison,
\emph{Objectivity}.}

Selected passages on mechanical objectivity and scientific images.
\parencite{dastongalison2007}

\textit{Type: History of science / visual culture}

\medskip

\textbf{Visual archive: selected motion studies by
Eadweard Muybridge and Étienne-Jules Marey.}

\textit{Type: Primary visual sources}

\begin{description}[leftmargin=1.15in,style=nextline]

\item[\textcolor{LangRed}{READ}]
What does ``mechanical objectivity'' promise? What practices of
seeing does it replace or reorganize?

\item[\textcolor{LangRed}{CONNECT}]
Compare a mathematical graph of motion, a sequence of photographs,
a table of measurements, and a numerical derivative. What form of
change does each make visible?

\item[\textcolor{LangRed}{QUESTION}]
Is an image evidence, measurement, interpretation, or model?
Can it be all four?

\end{description}

% ---------------------------------------------------------
\subsection*{Encounter 9 \hfill November 4}
\courselabel{Slope, extrema, and optimization}

\textbf{Michel Serres, \emph{The Birth of Physics}.}

Selected passage from ``Slope and Extrema,'' beginning around p.~52.

\textit{Type: Philosophy / history of science}

\begin{description}[leftmargin=1.15in,style=nextline]

\item[\textcolor{LangRed}{READ}]
How does Serres connect slope, deviation, equilibrium, and extrema
to a physical picture of the world?

\item[\textcolor{LangRed}{CONNECT}]
If a function is a landscape, its derivative tells us which way
the landscape slopes. An extremum is a place where that local
direction vanishes.

What changes when we turn this geometric fact into an algorithm:
repeatedly moving in the direction indicated by the derivative?

\item[\textcolor{LangRed}{QUESTION}]
Is optimization simply a mathematical technique, or does it encourage
a particular way of imagining a problem---as a landscape with states,
directions, extrema, and preferred paths?

\end{description}


% ---------------------------------------------------------
\subsection*{Encounter 8 \hfill November 9}
\courselabel{What does it mean for a machine to learn?}

\textbf{Frank Rosenblatt, ``The Perceptron: A Probabilistic Model for
Information Storage and Organization in the Brain.''}
Selected opening and model passages from the 1958 paper.

\textit{Type: Primary source / history of artificial intelligence}

\begin{description}[leftmargin=1.15in,style=nextline]
\item[\textcolor{LangRed}{READ}]
What does Rosenblatt claim the perceptron can model, and what does he mean by
learning or self-organization?
\item[\textcolor{LangRed}{CONNECT}]
In our perceptron model, identify the inputs, weights, weighted sum, decision
rule, error, and update. Exactly what changes when the machine ``learns''?
\item[\textcolor{LangRed}{QUESTION}]
Formulate the strongest mathematical argument you can for calling parameter
adjustment ``learning,'' and the strongest conceptual argument against doing so.
\end{description}

% ---------------------------------------------------------
\subsection*{Encounter 9 \hfill November 11}
\courselabel{The chain rule becomes a learning procedure}

\textbf{David E. Rumelhart, Geoffrey E. Hinton, and Ronald J. Williams,
``Learning Representations by Back-Propagating Errors.''}
\emph{Nature} 323 (1986), pp.~533--536.

\textit{Type: Primary technical source / artificial intelligence}

\begin{description}[leftmargin=1.15in,style=nextline]
\item[\textcolor{LangRed}{READ}]
What problem does backpropagation solve, and what do the authors claim hidden
units learn to represent?
\item[\textcolor{LangRed}{CONNECT}]
Annotate the paper using the mathematics of the course: Where is the function?
Where is composition? Where is the loss? Where is the derivative? Where is the
chain rule? Where is optimization?
\item[\textcolor{LangRed}{QUESTION}]
What has to happen conceptually before a sequence of derivative calculations
can be described as ``learning representations''?
\end{description}

% ---------------------------------------------------------
\subsection*{Encounter 10 \hfill November 16}
\courselabel{Flows and paths}

\textbf{Michel Serres, \emph{The Birth of Physics}.}
Working selection from ``Flows and Paths,'' beginning around p.~70.

\textit{Type: Philosophy / history of science}

\begin{description}[leftmargin=1.15in,style=nextline]
\item[\textcolor{LangRed}{READ}]
What distinction does Serres make between describing an object and describing
a flow, path, or process?
\item[\textcolor{LangRed}{CONNECT}]
A differential equation gives a rule for change rather than directly giving a
trajectory. How does this alter what it means to know a mathematical object?
\item[\textcolor{LangRed}{QUESTION}]
Does the language of flow illuminate differential equations, or become merely
metaphorical once we formalize the mathematics?
\end{description}

% ---------------------------------------------------------
\subsection*{Encounter 11 \hfill November 18}
\courselabel{Can a model be verified?}

\textbf{Naomi Oreskes, Kristin Shrader-Frechette, and Kenneth Belitz,
``Verification, Validation, and Confirmation of Numerical Models in the Earth
Sciences.''}
\emph{Science} 263 (1994), pp.~641--646.

\textit{Type: Philosophy / history of modeling}

\begin{description}[leftmargin=1.15in,style=nextline]
\item[\textcolor{LangRed}{READ}]
Why do the authors distinguish confirmation from verification or validation for
models of open natural systems?
\item[\textcolor{LangRed}{CONNECT}]
Return to an Euler-method simulation. Which parts come from mathematics, which
from empirical assumptions, which from numerical choices, and which from
initial conditions?
\item[\textcolor{LangRed}{QUESTION}]
If a model produces accurate predictions, what has actually been established?
What has not?
\end{description}

% ---------------------------------------------------------
\subsection*{Encounter 12 \hfill November 30}
\courselabel{Feedback, control, and intelligence}

\textbf{Norbert Wiener, \emph{The Human Use of Human Beings}.}
Short selection on cybernetics, feedback, communication, and purposive behavior.
\parencite{wiener1950}

\textit{Type: Primary source / cybernetics}

\medskip

\textbf{Andrew Pickering, \emph{The Cybernetic Brain}.}
Short historical selection on cybernetic practice.

\textit{Type: History / sociology of science}

\begin{description}[leftmargin=1.15in,style=nextline]
\item[\textcolor{LangRed}{READ}]
How does feedback permit Wiener and other cyberneticians to describe purposive
behavior without beginning from an internal mind or intention?
\item[\textcolor{LangRed}{CONNECT}]
Identify feedback mathematically in one of our dynamical systems. How do
equilibrium and stability change when the output of a system becomes part of
its future input?
\item[\textcolor{LangRed}{QUESTION}]
What is gained---and what may be lost---when intelligence, behavior, or purpose
is redescribed in terms of information and feedback?
\end{description}

% ---------------------------------------------------------
% ---------------------------------------------------------
\subsection*{Encounter 13 \hfill December 2}
\courselabel{When an equation becomes a field}

\textbf{Jean le Rond d'Alembert, Leonhard Euler, and Daniel Bernoulli,
``The Vibrating String and Its Partial Differential Equation,''
in \textcite{struik1986}, selection beginning p.~351.}

\textit{Type: Primary mathematical / physical source}

\begin{description}[leftmargin=1.15in,style=nextline]

\item[\textcolor{LangRed}{READ}]
What is being represented mathematically when the vibrating string is
described not by a single changing quantity but by a function of both
position and time?

\item[\textcolor{LangRed}{CONNECT}]
Compare

\[
y=y(t)
\]

with

\[
u=u(x,t).
\]

Why does the second description require partial derivatives? What
does each independent variable represent?

\item[\textcolor{LangRed}{QUESTION}]
Is the partial differential equation a description of the string,
a model of the string, or a new mathematical object with a life
independent of the physical system that produced it?

\end{description}



% ---------------------------------------------------------
\subsection*{Encounter 14 \hfill December 7}
\courselabel{Surfaces, dimensions, and classification}

\textbf{Gaspard Monge, ``The Two Curvatures of a Curved Surface,''
in \textcite{struik1986}, short selection beginning p.~413.}

\textit{Type: Primary mathematical source}

\medskip

\textbf{Geoffrey C. Bowker and Susan Leigh Star,
\emph{Sorting Things Out: Classification and Its Consequences}.}

Selected introduction. \parencite{bowkerstar1999}

\textit{Type: Science and technology studies}

\begin{description}[leftmargin=1.15in,style=nextline]

\item[\textcolor{LangRed}{READ}]
How does moving from a curve to a surface change the mathematical
problem? How do Bowker and Star describe the work performed by
systems of classification?

\item[\textcolor{LangRed}{CONNECT}]
For a curve

\[
y=f(x),
\]

local change can be described by a single derivative. For a surface

\[
z=f(x,y),
\]

change depends upon direction.

How does moving from curves to surfaces change the meaning of slope,
tangent, and curvature? How does this prepare us to think about
high-dimensional spaces in which objects are represented by many
coordinates?

\item[\textcolor{LangRed}{QUESTION}]
An embedding represents an object as a point in a high-dimensional
space. What decisions must occur before similarity, difference, and
classification can appear as geometric relationships?

\end{description}

% ---------------------------------------------------------
\subsection*{Encounter 17 \hfill December 9}
\courselabel{Images for machines}

\textbf{Trevor Paglen, ``Invisible Images (Your Pictures Are Looking at
You).''}

\textit{Type: Art / visual culture / critical technology studies}

\begin{description}[leftmargin=1.15in,style=nextline]
\item[\textcolor{LangRed}{READ}]
What changes, for Paglen, when images circulate primarily between machines
rather than between images and human viewers?
\item[\textcolor{LangRed}{CONNECT}]
Trace the transformations
\[
\text{image} \longrightarrow \text{pixels} \longrightarrow
\text{numbers} \longrightarrow \text{features} \longrightarrow
\text{vector / embedding} \longrightarrow \text{classification}.
\]
At which stages is calculus involved?
\item[\textcolor{LangRed}{QUESTION}]
If no human needs to see an image for the image to act in the world, should we
still understand it primarily as a visual object?
\end{description}

% ---------------------------------------------------------
\subsection*{Encounter 18 \hfill December 14}
\courselabel{Generative AI: Mathematics, infrastructure, and culture}

\textbf{Kate Crawford, \emph{Atlas of AI}.}
Selected concluding or synthetic passage. \parencite{crawford2021}

\textit{Type: Critical AI / sociology of technology}

\medskip

\textbf{Contemporary primary technical artifact.}
A short excerpt, diagram, or model description from work on diffusion or
generative models, to be selected.

\textit{Type: Primary technical source}

\begin{description}[leftmargin=1.15in,style=nextline]
\item[\textcolor{LangRed}{READ}]
What material, institutional, political, and labor systems does Crawford ask us
to see when we use the term ``artificial intelligence''?
\item[\textcolor{LangRed}{CONNECT}]
Our mathematical story may describe a generative system through probability,
gradients, fields, differential equations, or optimization. What does that
description explain extremely well? What does it leave outside the frame?
\item[\textcolor{LangRed}{QUESTION}]
What would an adequate account of generative AI have to contain in addition to
its mathematics?
\end{description}

% =========================================================
\section{Calculus Colloquium: Extended Reader}

The shared reader is intentionally small. The Calculus Colloquium expands it.
Students will select additional primary sources, historical scholarship,
theoretical texts, mathematical problems, visual artifacts, datasets, or
technical papers related to the current module.

Possible directions include Archimedes, Lucretius, Cavalieri, Newton, Leibniz,
Euler, Fourier, Marey, Muybridge, Turing, Shannon, McCulloch and Pitts, Hebb,
Rosenblatt, Wiener, Minsky and Papert, Rumelhart--Hinton--Williams, Hopfield,
convolutional networks, word embeddings, attention, transformers, GANs,
diffusion models, cybernetic art, Project Cybersyn, ImageNet, algorithmic
classification, scientific visualization, and histories of computation.

A colloquium contribution should not simply summarize its source. It should
make an argument about the relationship between the source and the mathematics
of the course.

%
% Citations assume references.bib is loaded in syllabus.tex via:
%     \addbibresource{references.bib}
%
% Page ranges below are working selections. Verify against the
% exact editions placed on Canvas / Course Reserves before release.
% =========================================================

\section{Reader}

The readings in this course are not supplementary background to the
mathematics. They are part of our investigation of what calculus is, how
mathematical knowledge develops, how models and measurements acquire authority,
and how ideas of change, optimization, representation, and intelligence have
been transformed across different historical settings.

The reader moves among three kinds of texts:

\begin{itemize}[leftmargin=*]
    \item \textbf{Primary sources}: mathematical, scientific, and technical
    texts written by historical actors as ideas were being developed;
    \item \textbf{History and sociology of science}: scholarship examining
    mathematics, models, objectivity, quantification, and scientific practice
    historically and socially;
    \item \textbf{Theory and cultural studies}: texts that help us question
    space, classification, representation, images, technology, and artificial
    intelligence.
\end{itemize}

Most reading assignments are deliberately short. Read slowly and bring the
text into contact with the mathematics. For each encounter, use the three
questions below:

\begin{description}[leftmargin=1.15in,style=nextline]
    \item[\textcolor{LangRed}{READ}] What is the author actually claiming,
    constructing, or trying to accomplish?
    \item[\textcolor{LangRed}{CONNECT}] What happens when we place the text
    beside the mathematics, computation, model, or image we are studying?
    \item[\textcolor{LangRed}{QUESTION}] Where should we challenge, test,
    complicate, or refuse the author's account?
\end{description}

% ---------------------------------------------------------
\subsection*{Encounter 1 \hfill August 26}
\courselabel{Genesis: What is calculus?}

\textbf{Otto Toeplitz, \emph{The Calculus: A Genetic Approach}.}
Introductory selection. \parencite{toeplitz1963}

\textit{Type: History of mathematics / pedagogy}

\begin{description}[leftmargin=1.15in,style=nextline]
\item[\textcolor{LangRed}{READ}]
What does Toeplitz mean by a ``genetic'' approach to mathematics?
\item[\textcolor{LangRed}{CONNECT}]
Compare the order in which mathematical ideas historically arise with the
order in which a modern calculus textbook usually presents them.
\item[\textcolor{LangRed}{QUESTION}]
Should the historical genesis of a mathematical idea determine how we learn
it? Where might historical order help, and where might it mislead?
\end{description}

% ---------------------------------------------------------
\subsection*{Encounter 2 \hfill September 2}
\courselabel{Archimedes: Measurement, exhaustion, and matter}

\textbf{Archimedes, \emph{Quadrature of the Parabola}.}
Short primary-source selection from the geometric argument.

\textit{Type: Primary mathematical source}

\medskip

\textbf{Michel Serres, \emph{The Birth of Physics}.}
``The Work of Archimedes,'' working selection pp.~32--36, with a short
selection from ``Archimedes on Thinking Deviation.''

\textit{Type: Philosophy / history of science}

\begin{description}[leftmargin=1.15in,style=nextline]
\item[\textcolor{LangRed}{READ}]
How does Archimedes reason about an area without possessing our modern
definition of the definite integral? What does Serres think is at stake in
Archimedean mathematics?
\item[\textcolor{LangRed}{CONNECT}]
Identify the finite approximations in Archimedes' argument. How are they like,
and unlike, the computational approximations we construct in Jupyter?
\item[\textcolor{LangRed}{QUESTION}]
Does describing Archimedes as someone who ``almost discovered calculus'' help
us understand his work, or erase the mathematical world in which he worked?
\end{description}

% ---------------------------------------------------------
% ---------------------------------------------------------
\subsection*{Encounter 3 \hfill September 14}
\courselabel{Indivisibles, the continuum, and nomad science}

\textbf{Bonaventura Cavalieri, ``Principle of Cavalieri'' and
``Integration,'' in \textcite{struik1986}, selections beginning
pp.~209 and 214.}

\textit{Type: Primary mathematical source}

\medskip

\textbf{Gilles Deleuze and Félix Guattari, ``1227: Treatise on
Nomadology---The War Machine,'' in \emph{A Thousand Plateaus}.}

Selected passages on nomad science and royal or State science.
\parencite{deleuzeguattari1987}

\textit{Type: Philosophy / cultural theory}

\begin{description}[leftmargin=1.15in,style=nextline]

\item[\textcolor{LangRed}{READ}]
How does Cavalieri reason with indivisibles? What distinction do
Deleuze and Guattari draw between nomad and royal forms of knowledge?

\item[\textcolor{LangRed}{CONNECT}]
Does Cavalieri's treatment of continuous magnitudes fit Deleuze and
Guattari's description of a mathematical practice concerned with
becoming, variation, and local operations rather than fixed forms?

\item[\textcolor{LangRed}{QUESTION}]
Does ``nomad science'' actually illuminate Cavalieri's mathematics,
or are we imposing a twentieth-century philosophical category on a
seventeenth-century practice?

\end{description}

% ---------------------------------------------------------
\subsection*{Encounter 4 \hfill September 16}
\courselabel{Smooth and striated space}

\textbf{Deleuze and Guattari, \emph{A Thousand Plateaus}.}
Short selected passage on smooth and striated space.

\textit{Type: Philosophy / cultural theory}

\begin{description}[leftmargin=1.15in,style=nextline]
\item[\textcolor{LangRed}{READ}]
How do smooth and striated spaces differ in the authors' account?
\item[\textcolor{LangRed}{CONNECT}]
A Riemann sum imposes a coordinate system, partitions an interval, samples a
function, assigns rectangles, and aggregates measurements. In what sense does
this ``striate'' a continuous object?
\item[\textcolor{LangRed}{QUESTION}]
What becomes mathematically possible because of the partition? What information
or experience of the continuous object might the partition suppress?
\end{description}

% ---------------------------------------------------------
\subsection*{Encounter 5 \hfill September 23}
\courselabel{Before the Fundamental Theorem was fundamental}

\textbf{Isaac Barrow, ``The Fundamental Theorem of the Calculus,''
in \textcite{struik1986}, selection beginning p.~253.}

\textit{Type: Primary mathematical source}

\begin{description}[leftmargin=1.15in,style=nextline]

\item[\textcolor{LangRed}{READ}]
What relationship does Barrow establish between problems involving
tangents and problems involving areas?

\item[\textcolor{LangRed}{CONNECT}]
Translate Barrow's geometric reasoning into our modern relationship
between accumulation and differentiation:

\[
F(x)=\int_a^x f(t)\,dt
\qquad\Longrightarrow\qquad
F'(x)=f(x).
\]

What does our notation make easier to see?

\item[\textcolor{LangRed}{QUESTION}]
What gets lost when we call Barrow's argument ``the Fundamental
Theorem of Calculus'' before the modern objects called derivatives
and integrals have been stabilized?

\end{description}

% ---------------------------------------------------------
\subsection*{Encounter 6 \hfill September 30}
\courselabel{Quantification: When inequality becomes a number}

\textbf{Theodore M. Porter, \emph{Trust in Numbers}.}
Working selection from ``Cultures of Objectivity'' and ``How Social Numbers
Are Made Valid.'' \parencite{porter1995}

\textit{Type: History / sociology of quantification}

\begin{description}[leftmargin=1.15in,style=nextline]
\item[\textcolor{LangRed}{READ}]
Why, for Porter, can quantitative procedures acquire authority precisely by
appearing to reduce personal judgment?
\item[\textcolor{LangRed}{CONNECT}]
Construct a Lorenz curve and calculate a Gini coefficient. Identify at least
three decisions required before ``inequality'' can be represented as an area
and then as a single number.
\item[\textcolor{LangRed}{QUESTION}]
Which decisions in the Gini calculation are mathematical and which are social?
Is that distinction as clear as it first appears?
\end{description}

% ---------------------------------------------------------


% ---------------------------------------------------------
% ---------------------------------------------------------
\subsection*{Encounter 7 \hfill October 14}
\courselabel{The differential calculus appears}

\textbf{Gottfried Wilhelm Leibniz, ``The First Publication of His
Differential Calculus,'' in \textcite{struik1986},
selection beginning p.~271.}

\textit{Type: Primary mathematical source}

\begin{description}[leftmargin=1.15in,style=nextline]

\item[\textcolor{LangRed}{READ}]
What mathematical objects and operations does Leibniz introduce?
How does his notation participate in the argument?

\item[\textcolor{LangRed}{CONNECT}]
Translate a portion of Leibniz's argument into contemporary notation.
Identify where we would now speak about a derivative, differential,
product rule, quotient rule, or tangent.

\item[\textcolor{LangRed}{QUESTION}]
Does translating Leibniz into modern notation help us understand him,
or does it make his mathematics look more like ours than it really was?

\end{description}

% ---------------------------------------------------------
\subsection*{Encounter 8 \hfill October 26}
\courselabel{Seeing change: Images and objectivity}

\textbf{Lorraine Daston and Peter Galison,
\emph{Objectivity}.}

Selected passages on mechanical objectivity and scientific images.
\parencite{dastongalison2007}

\textit{Type: History of science / visual culture}

\medskip

\textbf{Visual archive: selected motion studies by
Eadweard Muybridge and Étienne-Jules Marey.}

\textit{Type: Primary visual sources}

\begin{description}[leftmargin=1.15in,style=nextline]

\item[\textcolor{LangRed}{READ}]
What does ``mechanical objectivity'' promise? What practices of
seeing does it replace or reorganize?

\item[\textcolor{LangRed}{CONNECT}]
Compare a mathematical graph of motion, a sequence of photographs,
a table of measurements, and a numerical derivative. What form of
change does each make visible?

\item[\textcolor{LangRed}{QUESTION}]
Is an image evidence, measurement, interpretation, or model?
Can it be all four?

\end{description}

% ---------------------------------------------------------
\subsection*{Encounter 9 \hfill November 4}
\courselabel{Slope, extrema, and optimization}

\textbf{Michel Serres, \emph{The Birth of Physics}.}

Selected passage from ``Slope and Extrema,'' beginning around p.~52.

\textit{Type: Philosophy / history of science}

\begin{description}[leftmargin=1.15in,style=nextline]

\item[\textcolor{LangRed}{READ}]
How does Serres connect slope, deviation, equilibrium, and extrema
to a physical picture of the world?

\item[\textcolor{LangRed}{CONNECT}]
If a function is a landscape, its derivative tells us which way
the landscape slopes. An extremum is a place where that local
direction vanishes.

What changes when we turn this geometric fact into an algorithm:
repeatedly moving in the direction indicated by the derivative?

\item[\textcolor{LangRed}{QUESTION}]
Is optimization simply a mathematical technique, or does it encourage
a particular way of imagining a problem---as a landscape with states,
directions, extrema, and preferred paths?

\end{description}


% ---------------------------------------------------------
\subsection*{Encounter 8 \hfill November 9}
\courselabel{What does it mean for a machine to learn?}

\textbf{Frank Rosenblatt, ``The Perceptron: A Probabilistic Model for
Information Storage and Organization in the Brain.''}
Selected opening and model passages from the 1958 paper.

\textit{Type: Primary source / history of artificial intelligence}

\begin{description}[leftmargin=1.15in,style=nextline]
\item[\textcolor{LangRed}{READ}]
What does Rosenblatt claim the perceptron can model, and what does he mean by
learning or self-organization?
\item[\textcolor{LangRed}{CONNECT}]
In our perceptron model, identify the inputs, weights, weighted sum, decision
rule, error, and update. Exactly what changes when the machine ``learns''?
\item[\textcolor{LangRed}{QUESTION}]
Formulate the strongest mathematical argument you can for calling parameter
adjustment ``learning,'' and the strongest conceptual argument against doing so.
\end{description}

% ---------------------------------------------------------
\subsection*{Encounter 9 \hfill November 11}
\courselabel{The chain rule becomes a learning procedure}

\textbf{David E. Rumelhart, Geoffrey E. Hinton, and Ronald J. Williams,
``Learning Representations by Back-Propagating Errors.''}
\emph{Nature} 323 (1986), pp.~533--536.

\textit{Type: Primary technical source / artificial intelligence}

\begin{description}[leftmargin=1.15in,style=nextline]
\item[\textcolor{LangRed}{READ}]
What problem does backpropagation solve, and what do the authors claim hidden
units learn to represent?
\item[\textcolor{LangRed}{CONNECT}]
Annotate the paper using the mathematics of the course: Where is the function?
Where is composition? Where is the loss? Where is the derivative? Where is the
chain rule? Where is optimization?
\item[\textcolor{LangRed}{QUESTION}]
What has to happen conceptually before a sequence of derivative calculations
can be described as ``learning representations''?
\end{description}

% ---------------------------------------------------------
\subsection*{Encounter 10 \hfill November 16}
\courselabel{Flows and paths}

\textbf{Michel Serres, \emph{The Birth of Physics}.}
Working selection from ``Flows and Paths,'' beginning around p.~70.

\textit{Type: Philosophy / history of science}

\begin{description}[leftmargin=1.15in,style=nextline]
\item[\textcolor{LangRed}{READ}]
What distinction does Serres make between describing an object and describing
a flow, path, or process?
\item[\textcolor{LangRed}{CONNECT}]
A differential equation gives a rule for change rather than directly giving a
trajectory. How does this alter what it means to know a mathematical object?
\item[\textcolor{LangRed}{QUESTION}]
Does the language of flow illuminate differential equations, or become merely
metaphorical once we formalize the mathematics?
\end{description}

% ---------------------------------------------------------
\subsection*{Encounter 11 \hfill November 18}
\courselabel{Can a model be verified?}

\textbf{Naomi Oreskes, Kristin Shrader-Frechette, and Kenneth Belitz,
``Verification, Validation, and Confirmation of Numerical Models in the Earth
Sciences.''}
\emph{Science} 263 (1994), pp.~641--646.

\textit{Type: Philosophy / history of modeling}

\begin{description}[leftmargin=1.15in,style=nextline]
\item[\textcolor{LangRed}{READ}]
Why do the authors distinguish confirmation from verification or validation for
models of open natural systems?
\item[\textcolor{LangRed}{CONNECT}]
Return to an Euler-method simulation. Which parts come from mathematics, which
from empirical assumptions, which from numerical choices, and which from
initial conditions?
\item[\textcolor{LangRed}{QUESTION}]
If a model produces accurate predictions, what has actually been established?
What has not?
\end{description}

% ---------------------------------------------------------
\subsection*{Encounter 12 \hfill November 30}
\courselabel{Feedback, control, and intelligence}

\textbf{Norbert Wiener, \emph{The Human Use of Human Beings}.}
Short selection on cybernetics, feedback, communication, and purposive behavior.
\parencite{wiener1950}

\textit{Type: Primary source / cybernetics}

\medskip

\textbf{Andrew Pickering, \emph{The Cybernetic Brain}.}
Short historical selection on cybernetic practice.

\textit{Type: History / sociology of science}

\begin{description}[leftmargin=1.15in,style=nextline]
\item[\textcolor{LangRed}{READ}]
How does feedback permit Wiener and other cyberneticians to describe purposive
behavior without beginning from an internal mind or intention?
\item[\textcolor{LangRed}{CONNECT}]
Identify feedback mathematically in one of our dynamical systems. How do
equilibrium and stability change when the output of a system becomes part of
its future input?
\item[\textcolor{LangRed}{QUESTION}]
What is gained---and what may be lost---when intelligence, behavior, or purpose
is redescribed in terms of information and feedback?
\end{description}

% ---------------------------------------------------------
% ---------------------------------------------------------
\subsection*{Encounter 13 \hfill December 2}
\courselabel{When an equation becomes a field}

\textbf{Jean le Rond d'Alembert, Leonhard Euler, and Daniel Bernoulli,
``The Vibrating String and Its Partial Differential Equation,''
in \textcite{struik1986}, selection beginning p.~351.}

\textit{Type: Primary mathematical / physical source}

\begin{description}[leftmargin=1.15in,style=nextline]

\item[\textcolor{LangRed}{READ}]
What is being represented mathematically when the vibrating string is
described not by a single changing quantity but by a function of both
position and time?

\item[\textcolor{LangRed}{CONNECT}]
Compare

\[
y=y(t)
\]

with

\[
u=u(x,t).
\]

Why does the second description require partial derivatives? What
does each independent variable represent?

\item[\textcolor{LangRed}{QUESTION}]
Is the partial differential equation a description of the string,
a model of the string, or a new mathematical object with a life
independent of the physical system that produced it?

\end{description}



% ---------------------------------------------------------
\subsection*{Encounter 14 \hfill December 7}
\courselabel{Surfaces, dimensions, and classification}

\textbf{Gaspard Monge, ``The Two Curvatures of a Curved Surface,''
in \textcite{struik1986}, short selection beginning p.~413.}

\textit{Type: Primary mathematical source}

\medskip

\textbf{Geoffrey C. Bowker and Susan Leigh Star,
\emph{Sorting Things Out: Classification and Its Consequences}.}

Selected introduction. \parencite{bowkerstar1999}

\textit{Type: Science and technology studies}

\begin{description}[leftmargin=1.15in,style=nextline]

\item[\textcolor{LangRed}{READ}]
How does moving from a curve to a surface change the mathematical
problem? How do Bowker and Star describe the work performed by
systems of classification?

\item[\textcolor{LangRed}{CONNECT}]
For a curve

\[
y=f(x),
\]

local change can be described by a single derivative. For a surface

\[
z=f(x,y),
\]

change depends upon direction.

How does moving from curves to surfaces change the meaning of slope,
tangent, and curvature? How does this prepare us to think about
high-dimensional spaces in which objects are represented by many
coordinates?

\item[\textcolor{LangRed}{QUESTION}]
An embedding represents an object as a point in a high-dimensional
space. What decisions must occur before similarity, difference, and
classification can appear as geometric relationships?

\end{description}

% ---------------------------------------------------------
\subsection*{Encounter 17 \hfill December 9}
\courselabel{Images for machines}

\textbf{Trevor Paglen, ``Invisible Images (Your Pictures Are Looking at
You).''}

\textit{Type: Art / visual culture / critical technology studies}

\begin{description}[leftmargin=1.15in,style=nextline]
\item[\textcolor{LangRed}{READ}]
What changes, for Paglen, when images circulate primarily between machines
rather than between images and human viewers?
\item[\textcolor{LangRed}{CONNECT}]
Trace the transformations
\[
\text{image} \longrightarrow \text{pixels} \longrightarrow
\text{numbers} \longrightarrow \text{features} \longrightarrow
\text{vector / embedding} \longrightarrow \text{classification}.
\]
At which stages is calculus involved?
\item[\textcolor{LangRed}{QUESTION}]
If no human needs to see an image for the image to act in the world, should we
still understand it primarily as a visual object?
\end{description}

% ---------------------------------------------------------
\subsection*{Encounter 18 \hfill December 14}
\courselabel{Generative AI: Mathematics, infrastructure, and culture}

\textbf{Kate Crawford, \emph{Atlas of AI}.}
Selected concluding or synthetic passage. \parencite{crawford2021}

\textit{Type: Critical AI / sociology of technology}

\medskip

\textbf{Contemporary primary technical artifact.}
A short excerpt, diagram, or model description from work on diffusion or
generative models, to be selected.

\textit{Type: Primary technical source}

\begin{description}[leftmargin=1.15in,style=nextline]
\item[\textcolor{LangRed}{READ}]
What material, institutional, political, and labor systems does Crawford ask us
to see when we use the term ``artificial intelligence''?
\item[\textcolor{LangRed}{CONNECT}]
Our mathematical story may describe a generative system through probability,
gradients, fields, differential equations, or optimization. What does that
description explain extremely well? What does it leave outside the frame?
\item[\textcolor{LangRed}{QUESTION}]
What would an adequate account of generative AI have to contain in addition to
its mathematics?
\end{description}

% =========================================================
\section{Calculus Colloquium: Extended Reader}

The shared reader is intentionally small. The Calculus Colloquium expands it.
Students will select additional primary sources, historical scholarship,
theoretical texts, mathematical problems, visual artifacts, datasets, or
technical papers related to the current module.

Possible directions include Archimedes, Lucretius, Cavalieri, Newton, Leibniz,
Euler, Fourier, Marey, Muybridge, Turing, Shannon, McCulloch and Pitts, Hebb,
Rosenblatt, Wiener, Minsky and Papert, Rumelhart--Hinton--Williams, Hopfield,
convolutional networks, word embeddings, attention, transformers, GANs,
diffusion models, cybernetic art, Project Cybersyn, ImageNet, algorithmic
classification, scientific visualization, and histories of computation.

A colloquium contribution should not simply summarize its source. It should
make an argument about the relationship between the source and the mathematics
of the course.

%
% Citations assume references.bib is loaded in syllabus.tex via:
%     \addbibresource{references.bib}
%
% Page ranges below are working selections. Verify against the
% exact editions placed on Canvas / Course Reserves before release.
% =========================================================

\section{Reader}

The readings in this course are not supplementary background to the
mathematics. They are part of our investigation of what calculus is, how
mathematical knowledge develops, how models and measurements acquire authority,
and how ideas of change, optimization, representation, and intelligence have
been transformed across different historical settings.

The reader moves among three kinds of texts:

\begin{itemize}[leftmargin=*]
    \item \textbf{Primary sources}: mathematical, scientific, and technical
    texts written by historical actors as ideas were being developed;
    \item \textbf{History and sociology of science}: scholarship examining
    mathematics, models, objectivity, quantification, and scientific practice
    historically and socially;
    \item \textbf{Theory and cultural studies}: texts that help us question
    space, classification, representation, images, technology, and artificial
    intelligence.
\end{itemize}

Most reading assignments are deliberately short. Read slowly and bring the
text into contact with the mathematics. For each encounter, use the three
questions below:

\begin{description}[leftmargin=1.15in,style=nextline]
    \item[\textcolor{LangRed}{READ}] What is the author actually claiming,
    constructing, or trying to accomplish?
    \item[\textcolor{LangRed}{CONNECT}] What happens when we place the text
    beside the mathematics, computation, model, or image we are studying?
    \item[\textcolor{LangRed}{QUESTION}] Where should we challenge, test,
    complicate, or refuse the author's account?
\end{description}

% ---------------------------------------------------------
\subsection*{Encounter 1 \hfill August 26}
\courselabel{Genesis: What is calculus?}

\textbf{Otto Toeplitz, \emph{The Calculus: A Genetic Approach}.}
Introductory selection. \parencite{toeplitz1963}

\textit{Type: History of mathematics / pedagogy}

\begin{description}[leftmargin=1.15in,style=nextline]
\item[\textcolor{LangRed}{READ}]
What does Toeplitz mean by a ``genetic'' approach to mathematics?
\item[\textcolor{LangRed}{CONNECT}]
Compare the order in which mathematical ideas historically arise with the
order in which a modern calculus textbook usually presents them.
\item[\textcolor{LangRed}{QUESTION}]
Should the historical genesis of a mathematical idea determine how we learn
it? Where might historical order help, and where might it mislead?
\end{description}

% ---------------------------------------------------------
\subsection*{Encounter 2 \hfill September 2}
\courselabel{Archimedes: Measurement, exhaustion, and matter}

\textbf{Archimedes, \emph{Quadrature of the Parabola}.}
Short primary-source selection from the geometric argument.

\textit{Type: Primary mathematical source}

\medskip

\textbf{Michel Serres, \emph{The Birth of Physics}.}
``The Work of Archimedes,'' working selection pp.~32--36, with a short
selection from ``Archimedes on Thinking Deviation.''

\textit{Type: Philosophy / history of science}

\begin{description}[leftmargin=1.15in,style=nextline]
\item[\textcolor{LangRed}{READ}]
How does Archimedes reason about an area without possessing our modern
definition of the definite integral? What does Serres think is at stake in
Archimedean mathematics?
\item[\textcolor{LangRed}{CONNECT}]
Identify the finite approximations in Archimedes' argument. How are they like,
and unlike, the computational approximations we construct in Jupyter?
\item[\textcolor{LangRed}{QUESTION}]
Does describing Archimedes as someone who ``almost discovered calculus'' help
us understand his work, or erase the mathematical world in which he worked?
\end{description}

% ---------------------------------------------------------
% ---------------------------------------------------------
\subsection*{Encounter 3 \hfill September 14}
\courselabel{Indivisibles, the continuum, and nomad science}

\textbf{Bonaventura Cavalieri, ``Principle of Cavalieri'' and
``Integration,'' in \textcite{struik1986}, selections beginning
pp.~209 and 214.}

\textit{Type: Primary mathematical source}

\medskip

\textbf{Gilles Deleuze and Félix Guattari, ``1227: Treatise on
Nomadology---The War Machine,'' in \emph{A Thousand Plateaus}.}

Selected passages on nomad science and royal or State science.
\parencite{deleuzeguattari1987}

\textit{Type: Philosophy / cultural theory}

\begin{description}[leftmargin=1.15in,style=nextline]

\item[\textcolor{LangRed}{READ}]
How does Cavalieri reason with indivisibles? What distinction do
Deleuze and Guattari draw between nomad and royal forms of knowledge?

\item[\textcolor{LangRed}{CONNECT}]
Does Cavalieri's treatment of continuous magnitudes fit Deleuze and
Guattari's description of a mathematical practice concerned with
becoming, variation, and local operations rather than fixed forms?

\item[\textcolor{LangRed}{QUESTION}]
Does ``nomad science'' actually illuminate Cavalieri's mathematics,
or are we imposing a twentieth-century philosophical category on a
seventeenth-century practice?

\end{description}

% ---------------------------------------------------------
\subsection*{Encounter 4 \hfill September 16}
\courselabel{Smooth and striated space}

\textbf{Deleuze and Guattari, \emph{A Thousand Plateaus}.}
Short selected passage on smooth and striated space.

\textit{Type: Philosophy / cultural theory}

\begin{description}[leftmargin=1.15in,style=nextline]
\item[\textcolor{LangRed}{READ}]
How do smooth and striated spaces differ in the authors' account?
\item[\textcolor{LangRed}{CONNECT}]
A Riemann sum imposes a coordinate system, partitions an interval, samples a
function, assigns rectangles, and aggregates measurements. In what sense does
this ``striate'' a continuous object?
\item[\textcolor{LangRed}{QUESTION}]
What becomes mathematically possible because of the partition? What information
or experience of the continuous object might the partition suppress?
\end{description}

% ---------------------------------------------------------
\subsection*{Encounter 5 \hfill September 23}
\courselabel{Before the Fundamental Theorem was fundamental}

\textbf{Isaac Barrow, ``The Fundamental Theorem of the Calculus,''
in \textcite{struik1986}, selection beginning p.~253.}

\textit{Type: Primary mathematical source}

\begin{description}[leftmargin=1.15in,style=nextline]

\item[\textcolor{LangRed}{READ}]
What relationship does Barrow establish between problems involving
tangents and problems involving areas?

\item[\textcolor{LangRed}{CONNECT}]
Translate Barrow's geometric reasoning into our modern relationship
between accumulation and differentiation:

\[
F(x)=\int_a^x f(t)\,dt
\qquad\Longrightarrow\qquad
F'(x)=f(x).
\]

What does our notation make easier to see?

\item[\textcolor{LangRed}{QUESTION}]
What gets lost when we call Barrow's argument ``the Fundamental
Theorem of Calculus'' before the modern objects called derivatives
and integrals have been stabilized?

\end{description}

% ---------------------------------------------------------
\subsection*{Encounter 6 \hfill September 30}
\courselabel{Quantification: When inequality becomes a number}

\textbf{Theodore M. Porter, \emph{Trust in Numbers}.}
Working selection from ``Cultures of Objectivity'' and ``How Social Numbers
Are Made Valid.'' \parencite{porter1995}

\textit{Type: History / sociology of quantification}

\begin{description}[leftmargin=1.15in,style=nextline]
\item[\textcolor{LangRed}{READ}]
Why, for Porter, can quantitative procedures acquire authority precisely by
appearing to reduce personal judgment?
\item[\textcolor{LangRed}{CONNECT}]
Construct a Lorenz curve and calculate a Gini coefficient. Identify at least
three decisions required before ``inequality'' can be represented as an area
and then as a single number.
\item[\textcolor{LangRed}{QUESTION}]
Which decisions in the Gini calculation are mathematical and which are social?
Is that distinction as clear as it first appears?
\end{description}

% ---------------------------------------------------------


% ---------------------------------------------------------
% ---------------------------------------------------------
\subsection*{Encounter 7 \hfill October 14}
\courselabel{The differential calculus appears}

\textbf{Gottfried Wilhelm Leibniz, ``The First Publication of His
Differential Calculus,'' in \textcite{struik1986},
selection beginning p.~271.}

\textit{Type: Primary mathematical source}

\begin{description}[leftmargin=1.15in,style=nextline]

\item[\textcolor{LangRed}{READ}]
What mathematical objects and operations does Leibniz introduce?
How does his notation participate in the argument?

\item[\textcolor{LangRed}{CONNECT}]
Translate a portion of Leibniz's argument into contemporary notation.
Identify where we would now speak about a derivative, differential,
product rule, quotient rule, or tangent.

\item[\textcolor{LangRed}{QUESTION}]
Does translating Leibniz into modern notation help us understand him,
or does it make his mathematics look more like ours than it really was?

\end{description}

% ---------------------------------------------------------
\subsection*{Encounter 8 \hfill October 26}
\courselabel{Seeing change: Images and objectivity}

\textbf{Lorraine Daston and Peter Galison,
\emph{Objectivity}.}

Selected passages on mechanical objectivity and scientific images.
\parencite{dastongalison2007}

\textit{Type: History of science / visual culture}

\medskip

\textbf{Visual archive: selected motion studies by
Eadweard Muybridge and Étienne-Jules Marey.}

\textit{Type: Primary visual sources}

\begin{description}[leftmargin=1.15in,style=nextline]

\item[\textcolor{LangRed}{READ}]
What does ``mechanical objectivity'' promise? What practices of
seeing does it replace or reorganize?

\item[\textcolor{LangRed}{CONNECT}]
Compare a mathematical graph of motion, a sequence of photographs,
a table of measurements, and a numerical derivative. What form of
change does each make visible?

\item[\textcolor{LangRed}{QUESTION}]
Is an image evidence, measurement, interpretation, or model?
Can it be all four?

\end{description}

% ---------------------------------------------------------
\subsection*{Encounter 9 \hfill November 4}
\courselabel{Slope, extrema, and optimization}

\textbf{Michel Serres, \emph{The Birth of Physics}.}

Selected passage from ``Slope and Extrema,'' beginning around p.~52.

\textit{Type: Philosophy / history of science}

\begin{description}[leftmargin=1.15in,style=nextline]

\item[\textcolor{LangRed}{READ}]
How does Serres connect slope, deviation, equilibrium, and extrema
to a physical picture of the world?

\item[\textcolor{LangRed}{CONNECT}]
If a function is a landscape, its derivative tells us which way
the landscape slopes. An extremum is a place where that local
direction vanishes.

What changes when we turn this geometric fact into an algorithm:
repeatedly moving in the direction indicated by the derivative?

\item[\textcolor{LangRed}{QUESTION}]
Is optimization simply a mathematical technique, or does it encourage
a particular way of imagining a problem---as a landscape with states,
directions, extrema, and preferred paths?

\end{description}


% ---------------------------------------------------------
\subsection*{Encounter 8 \hfill November 9}
\courselabel{What does it mean for a machine to learn?}

\textbf{Frank Rosenblatt, ``The Perceptron: A Probabilistic Model for
Information Storage and Organization in the Brain.''}
Selected opening and model passages from the 1958 paper.

\textit{Type: Primary source / history of artificial intelligence}

\begin{description}[leftmargin=1.15in,style=nextline]
\item[\textcolor{LangRed}{READ}]
What does Rosenblatt claim the perceptron can model, and what does he mean by
learning or self-organization?
\item[\textcolor{LangRed}{CONNECT}]
In our perceptron model, identify the inputs, weights, weighted sum, decision
rule, error, and update. Exactly what changes when the machine ``learns''?
\item[\textcolor{LangRed}{QUESTION}]
Formulate the strongest mathematical argument you can for calling parameter
adjustment ``learning,'' and the strongest conceptual argument against doing so.
\end{description}

% ---------------------------------------------------------
\subsection*{Encounter 9 \hfill November 11}
\courselabel{The chain rule becomes a learning procedure}

\textbf{David E. Rumelhart, Geoffrey E. Hinton, and Ronald J. Williams,
``Learning Representations by Back-Propagating Errors.''}
\emph{Nature} 323 (1986), pp.~533--536.

\textit{Type: Primary technical source / artificial intelligence}

\begin{description}[leftmargin=1.15in,style=nextline]
\item[\textcolor{LangRed}{READ}]
What problem does backpropagation solve, and what do the authors claim hidden
units learn to represent?
\item[\textcolor{LangRed}{CONNECT}]
Annotate the paper using the mathematics of the course: Where is the function?
Where is composition? Where is the loss? Where is the derivative? Where is the
chain rule? Where is optimization?
\item[\textcolor{LangRed}{QUESTION}]
What has to happen conceptually before a sequence of derivative calculations
can be described as ``learning representations''?
\end{description}

% ---------------------------------------------------------
\subsection*{Encounter 10 \hfill November 16}
\courselabel{Flows and paths}

\textbf{Michel Serres, \emph{The Birth of Physics}.}
Working selection from ``Flows and Paths,'' beginning around p.~70.

\textit{Type: Philosophy / history of science}

\begin{description}[leftmargin=1.15in,style=nextline]
\item[\textcolor{LangRed}{READ}]
What distinction does Serres make between describing an object and describing
a flow, path, or process?
\item[\textcolor{LangRed}{CONNECT}]
A differential equation gives a rule for change rather than directly giving a
trajectory. How does this alter what it means to know a mathematical object?
\item[\textcolor{LangRed}{QUESTION}]
Does the language of flow illuminate differential equations, or become merely
metaphorical once we formalize the mathematics?
\end{description}

% ---------------------------------------------------------
\subsection*{Encounter 11 \hfill November 18}
\courselabel{Can a model be verified?}

\textbf{Naomi Oreskes, Kristin Shrader-Frechette, and Kenneth Belitz,
``Verification, Validation, and Confirmation of Numerical Models in the Earth
Sciences.''}
\emph{Science} 263 (1994), pp.~641--646.

\textit{Type: Philosophy / history of modeling}

\begin{description}[leftmargin=1.15in,style=nextline]
\item[\textcolor{LangRed}{READ}]
Why do the authors distinguish confirmation from verification or validation for
models of open natural systems?
\item[\textcolor{LangRed}{CONNECT}]
Return to an Euler-method simulation. Which parts come from mathematics, which
from empirical assumptions, which from numerical choices, and which from
initial conditions?
\item[\textcolor{LangRed}{QUESTION}]
If a model produces accurate predictions, what has actually been established?
What has not?
\end{description}

% ---------------------------------------------------------
\subsection*{Encounter 12 \hfill November 30}
\courselabel{Feedback, control, and intelligence}

\textbf{Norbert Wiener, \emph{The Human Use of Human Beings}.}
Short selection on cybernetics, feedback, communication, and purposive behavior.
\parencite{wiener1950}

\textit{Type: Primary source / cybernetics}

\medskip

\textbf{Andrew Pickering, \emph{The Cybernetic Brain}.}
Short historical selection on cybernetic practice.

\textit{Type: History / sociology of science}

\begin{description}[leftmargin=1.15in,style=nextline]
\item[\textcolor{LangRed}{READ}]
How does feedback permit Wiener and other cyberneticians to describe purposive
behavior without beginning from an internal mind or intention?
\item[\textcolor{LangRed}{CONNECT}]
Identify feedback mathematically in one of our dynamical systems. How do
equilibrium and stability change when the output of a system becomes part of
its future input?
\item[\textcolor{LangRed}{QUESTION}]
What is gained---and what may be lost---when intelligence, behavior, or purpose
is redescribed in terms of information and feedback?
\end{description}

% ---------------------------------------------------------
% ---------------------------------------------------------
\subsection*{Encounter 13 \hfill December 2}
\courselabel{When an equation becomes a field}

\textbf{Jean le Rond d'Alembert, Leonhard Euler, and Daniel Bernoulli,
``The Vibrating String and Its Partial Differential Equation,''
in \textcite{struik1986}, selection beginning p.~351.}

\textit{Type: Primary mathematical / physical source}

\begin{description}[leftmargin=1.15in,style=nextline]

\item[\textcolor{LangRed}{READ}]
What is being represented mathematically when the vibrating string is
described not by a single changing quantity but by a function of both
position and time?

\item[\textcolor{LangRed}{CONNECT}]
Compare

\[
y=y(t)
\]

with

\[
u=u(x,t).
\]

Why does the second description require partial derivatives? What
does each independent variable represent?

\item[\textcolor{LangRed}{QUESTION}]
Is the partial differential equation a description of the string,
a model of the string, or a new mathematical object with a life
independent of the physical system that produced it?

\end{description}



% ---------------------------------------------------------
\subsection*{Encounter 14 \hfill December 7}
\courselabel{Surfaces, dimensions, and classification}

\textbf{Gaspard Monge, ``The Two Curvatures of a Curved Surface,''
in \textcite{struik1986}, short selection beginning p.~413.}

\textit{Type: Primary mathematical source}

\medskip

\textbf{Geoffrey C. Bowker and Susan Leigh Star,
\emph{Sorting Things Out: Classification and Its Consequences}.}

Selected introduction. \parencite{bowkerstar1999}

\textit{Type: Science and technology studies}

\begin{description}[leftmargin=1.15in,style=nextline]

\item[\textcolor{LangRed}{READ}]
How does moving from a curve to a surface change the mathematical
problem? How do Bowker and Star describe the work performed by
systems of classification?

\item[\textcolor{LangRed}{CONNECT}]
For a curve

\[
y=f(x),
\]

local change can be described by a single derivative. For a surface

\[
z=f(x,y),
\]

change depends upon direction.

How does moving from curves to surfaces change the meaning of slope,
tangent, and curvature? How does this prepare us to think about
high-dimensional spaces in which objects are represented by many
coordinates?

\item[\textcolor{LangRed}{QUESTION}]
An embedding represents an object as a point in a high-dimensional
space. What decisions must occur before similarity, difference, and
classification can appear as geometric relationships?

\end{description}

% ---------------------------------------------------------
\subsection*{Encounter 17 \hfill December 9}
\courselabel{Images for machines}

\textbf{Trevor Paglen, ``Invisible Images (Your Pictures Are Looking at
You).''}

\textit{Type: Art / visual culture / critical technology studies}

\begin{description}[leftmargin=1.15in,style=nextline]
\item[\textcolor{LangRed}{READ}]
What changes, for Paglen, when images circulate primarily between machines
rather than between images and human viewers?
\item[\textcolor{LangRed}{CONNECT}]
Trace the transformations
\[
\text{image} \longrightarrow \text{pixels} \longrightarrow
\text{numbers} \longrightarrow \text{features} \longrightarrow
\text{vector / embedding} \longrightarrow \text{classification}.
\]
At which stages is calculus involved?
\item[\textcolor{LangRed}{QUESTION}]
If no human needs to see an image for the image to act in the world, should we
still understand it primarily as a visual object?
\end{description}

% ---------------------------------------------------------
\subsection*{Encounter 18 \hfill December 14}
\courselabel{Generative AI: Mathematics, infrastructure, and culture}

\textbf{Kate Crawford, \emph{Atlas of AI}.}
Selected concluding or synthetic passage. \parencite{crawford2021}

\textit{Type: Critical AI / sociology of technology}

\medskip

\textbf{Contemporary primary technical artifact.}
A short excerpt, diagram, or model description from work on diffusion or
generative models, to be selected.

\textit{Type: Primary technical source}

\begin{description}[leftmargin=1.15in,style=nextline]
\item[\textcolor{LangRed}{READ}]
What material, institutional, political, and labor systems does Crawford ask us
to see when we use the term ``artificial intelligence''?
\item[\textcolor{LangRed}{CONNECT}]
Our mathematical story may describe a generative system through probability,
gradients, fields, differential equations, or optimization. What does that
description explain extremely well? What does it leave outside the frame?
\item[\textcolor{LangRed}{QUESTION}]
What would an adequate account of generative AI have to contain in addition to
its mathematics?
\end{description}

% =========================================================
\section{Calculus Colloquium: Extended Reader}

The shared reader is intentionally small. The Calculus Colloquium expands it.
Students will select additional primary sources, historical scholarship,
theoretical texts, mathematical problems, visual artifacts, datasets, or
technical papers related to the current module.

Possible directions include Archimedes, Lucretius, Cavalieri, Newton, Leibniz,
Euler, Fourier, Marey, Muybridge, Turing, Shannon, McCulloch and Pitts, Hebb,
Rosenblatt, Wiener, Minsky and Papert, Rumelhart--Hinton--Williams, Hopfield,
convolutional networks, word embeddings, attention, transformers, GANs,
diffusion models, cybernetic art, Project Cybersyn, ImageNet, algorithmic
classification, scientific visualization, and histories of computation.

A colloquium contribution should not simply summarize its source. It should
make an argument about the relationship between the source and the mathematics
of the course.
